\subsection{SoC interface}
\paragraph{}
The soc\_if.v module is responsible for AXI slave decoding and interrupt control. It is a Wishbone compliant memory mapped module.  
The block diagram of this module is shown in  Figure \ref{fig:socif_fig1}. 


\begin{figure}[h]\centering % Using \begin{figure*} makes the figure take up the entire width of the page
	\includegraphics[width=\linewidth]{socif_fig1}
	\caption{SoC interface}
	\label{fig:socif_fig1}
\end{figure}

\paragraph{Wishbone Slave Interface}
The soc\_if module is connected to the system Wishbone bus via this Wishbone slave interface. \\
The port-list of the Wishbone slave interface is shown Table \ref{tab:socif_northbound}.

\begin{table}[hbt]
\caption{Wishbone Slave Interface Ports}
%	\centering
\begin{tabular}{|l|l|l|l|}
\toprule
%		\multicolumn{4}{c}{Name} \\
%		\cmidrule(r){1-4}
Port  & Width &Direction & Description  \\
\midrule
	\hline WB\_CLK\_I & & input & The clock input [CLK\_I] coordinates all activities \\ & & & for the internal logic within the
WISHBONE interface.\\
	\hline WB\_RST\_I & & input & The reset input [RST\_I] forces the WISHBONE interface to restart\\
	\hline WBS\_STB\_I&  & input & The strobe input [STB\_I], when asserted, \\  & & &  indicates that the SLAVE is selected\\
	\hline WBS\_CYC\_I&  & input & The cycle input [CYC\_I], when asserted,  \\  & & & indicates that a valid bus cycle is in progress\\
	\hline WBS\_WE\_I & & input & The write enable input [WE\_I] indicates whether  \\  & & & the current local bus cycle is  a READ or
WRITE cycle\\
	\hline WBS\_SEL\_I&  [WB\_DATA\_WIDTH/8-1:0] & input & The select input array [SEL\_I()] indicates where valid data is place\\
	\hline WBS\_DAT\_I& [WB\_DATA\_WIDTH-1:0] & input & The data input array [DAT\_I()] is used to pass binary data\\
	\hline WBS\_ADR\_I& [WB\_ADDR\_WIDTH-1:0] & input & The address input array [ADR\_I()] is used to pass a binary address \\
	\hline WBS\_ACK\_O&  & output & The acknowledge output [ACK\_O], when asserted, indicates  \\  & & &  the termination of a normal
bus cycle\\
	\hline WBS\_DAT\_O&  [WB\_DATA\_WIDTH-1:0]& output & The data output array [DAT\_O()] is used to pass binary data\\
\bottomrule
\end{tabular}
\label{tab:socif_northbound}
\end{table}

\paragraph{AXI-lite interfaces} 
The soc\_if module generates S\_COUNT AXI-lite memory mapped interfaces
starting at 0x80100000. These AXI-lite interfaces are connected to the Ibuf,
conv2d CNN layers and the Obuf modules as shown in the Figure \ref{fig:figure2}. All the
configurations of the internal modules are done via these AXI-lite
interfaces.\\
The  base addresses for memory map CSR and axi-lite slave regions are define in table \ref{tab:slv_addr_space}

\begin{table}[hbt!]
\caption{AXI-lite slave region base addresses}
\centering
\begin{tabular}{|l|l|}
\toprule
%		\multicolumn{2}{c}{Name} \\
%		\cmidrule(r){1-2}
Base Address  & Slave   \\
\midrule
	\hline0x80000000 & Memory Mapped  \\
	  &   \\
	 &  \\
	\hline 0x80100000& AXI-lite Slave 0 \\
	  &   \\
	&  \\
	\hline 0x80200000 & AXI-lite Slave 1 \\
	  &   \\
	 &  \\
	\hline 0x80300000& AXI-lite Slave 2 \\
	  &   \\
	 &  \\
	\hline 0x80400000 & AXI-lite Slave 3 \\
	  &   \\
	 &  \\
	\hline 0x80500000& AXI-lite Slave 4 \\
	 &   \\
	 &  \\
	\hline .... & .... \\
	\hline 0x80N00000 & AXI-lite Slave N\\
	  &   \\
	 &  \\
\bottomrule
\end{tabular}
\label{tab:slv_addr_space}
\end{table}

The port-list of these AXI-lite interfaces is shown in Table \ref{tab:socif_southbound}.  

\begin{table}[hbt]

\caption{AXI-Lite Interface Ports}
%	\centering
\begin{tabular}{|l|l|l|l|}
\toprule
%		\multicolumn{4}{c}{Name} \\
%		\cmidrule(r){1-4}
Port  & Width & Direction & Description  \\
\midrule
	\hline \multicolumn{4}{l}{Global Signals}\\
	\hline M\_AXI\_ACLK & & output & Clock signal. \\
	\hline M\_AXI\_ARESETN&  & output & Reset signal, active LOW\\
	\hline \multicolumn{4}{l}{AXI-lite write address channel}\\
	\hline M\_AXI\_AWVALID&  [S\_COUNT-1:0] & output & Write address valid \\
	\hline M\_AXI\_AWREADY&  [S\_COUNT-1:0] & input & Write address ready\\
	\hline M\_AXI\_AWADDR&  [AXI\_ADDR\_WIDTH-1:0] & output &Write address \\
	\hline M\_AXI\_AWPROT&  [2:0] & output & Protection type\\
	\hline \multicolumn{4}{l}{AXI-lite write data channel}\\
	\hline M\_AXI\_WVALID&  [S\_COUNT-1:0] & output & Write valid \\
	\hline M\_AXI\_WREADY&  [S\_COUNT-1:0] & input & Write ready. \\
	\hline M\_AXI\_WDATA&  [AXI\_DATA\_WIDTH-1:0] & output & Write data\\
	\hline M\_AXI\_WSTRB&   & output & Write strobes\\
	\hline \multicolumn{4}{l}{AXI-lite write response channel}\\
	\hline M\_AXI\_BVALID&  [S\_COUNT-1:0]& input & Write response valid\\
	\hline M\_AXI\_BREADY&  [S\_COUNT-1:0] & output & Response ready\\
	\hline M\_AXI\_BRESP&  [S\_COUNT*2-1:0] & input & Write response \\
	\hline \multicolumn{4}{l}{AXI-lite read address channel}\\
	\hline M\_AXI\_ARVALID&  [S\_COUNT-1:0]& output & Read address valid\\
	\hline M\_AXI\_ARREADY&  [S\_COUNT-1:0] & input & Read address ready\\
	\hline M\_AXI\_ARADDR&  [AXI\_ADDR\_WIDTH-1:0] & output & Read address\\
	\hline M\_AXI\_ARPROT& [2:0] & Output & Protection type\\
	\hline \multicolumn{4}{l}{AXI-lite read response channel}\\
	\hline M\_AXI\_RVALID&  [S\_COUNT-1:0] & 0utput & Read valid\\
	\hline M\_AXI\_RREADY&  [S\_COUNT-1:0]& input & Read ready \\
	\hline M\_AXI\_RDATA&  [S\_COUNT*AXI\_DATA\_WIDTH-1:0] & output & Read data\\
	\hline M\_AXI\_RRESP&  [S\_COUNT*2-1:0]  & input & Read response\\

\bottomrule
\end{tabular}
\label{tab:socif_southbound}
\end{table}


\paragraph{Wishbone to AXI-lite translator and decoder } 
The Wishbone to AXI-lite translator and decoder accepts the Wishbone transactions and decodes on the basis of the address to transit the transactions to the specific AXI-lite interface of the soc\_if module. \\

\paragraph{Interrupt Control}  
The interrupts from the Ibuf and Obuf are captured and registored in a memory
mapped register with in the soc\_if. This register can be can access by the
system through the Wishbone interconnect.  An output interrupt signal is also
generated based on the status of the interrupt register.
TODO: need to add register definition table with interrupt mask and clear bits
  
The interrupt ports are shown in Table \ref{tab:socif_intr}.  

\begin{table}[hbt]
\caption{Interrupt ports}
%	\centering
\begin{tabular}{|l|l|l|}
\toprule
%		\multicolumn{3}{c}{Name} \\
%		\cmidrule(r){1-3}
Port  & Direction & Description  \\
\midrule
	\hline ib\_intrpt & input & IBuf interrupt \\
	\hline ob\_intrpt & input & oBuf interrupt\\
	\hline soc\_if\_intrpt & Output & soc\_if\_intrpt\\

\bottomrule
\end{tabular}
\label{tab:socif_intr}
\end{table}
